\section{Introduction}
Visual speech recognition (VSR) recovers linguistic content from visible mouth
movements. Large-scale supervised learning has substantially improved benchmark
accuracy \cite{shillingford2019large,ma2023autoavsr}, yet strong performance on a
training corpus does not ensure reliable transcription of a different video
collection \cite{djilali2024wildvsr}. For an existing recognizer, a practical
question is how to reuse labelled source videos before deployment elsewhere.
We ask whether changing the temporal scope of supervision during source-only
continuation improves complete destination transcripts, and what it gives up
on the source corpus.

Complete utterances supply extensive visual and linguistic context. Short views
contain fewer neighboring frames and shorter target sequences, changing both
the available evidence and the optimization problem. A temporal cut can also
remove visual evidence for a boundary word retained in the target. These changes
make short-view supervision a testable intervention, rather than an assumed
improvement over complete clips.

Our initial formulation, fragment-mixed continuation (FM), alternates complete
source utterances with fixed-duration, transcript-labelled fragments of the same
videos. Fragment targets contain the transcript units overlapping each interval.
Only the training presentations change: the recognizer keeps its architecture
and ordinary whole-utterance decoding. No fragment selection, timestamp estimation
or prediction merging is required at inference, and no additional source videos
are introduced.

Full-clip continuation (FC) provides the comparison under the same initialization,
source-clip order and 6,000 optimizer updates. A common decoder compares fitted
states; development-selected decoders additionally describe calibrated systems.
Equal numbers of updates and presentations do not imply equal frame or token
exposure. We measure these exposures and retain the native objective, treating
the recipe as a complete intervention under an update budget rather than
attributing its effect to one regularization mechanism.

Two corpus pairs expose different outcomes. In VLR-to-ViCocktail transfer,
FM reduces common-decoder WER by 3.41 and 3.17 percentage points across two
seeds. In LRS2-to-WildVSR transfer, it gives no clear aggregate gain. The seed-17
source assessment shows that FM improves over each initializer but retains less
of FC's source improvement in both settings. The Vietnamese transfer gain thus
accompanies a source cost. Differences in initialization, training coverage,
repetition and calibration prevent attributing the contrast between corpus pairs
to language alone.

Temporal cropping and mixed-length training are established in lip reading
\cite{martinez2020lipreading,prajwal2022subword}. A separate matched comparison
therefore includes conventional whole-word mixing, which improves transfer and
outperforms fixed-time mixing by 0.459 WER points. A paired token-normalized loss
control retains the fixed-time gain. The whole-word short-only control further
tests the contribution of scheduled full-clip presentations, using the same
source order and unchanged whole-utterance assessment. Together, these controls
test how temporal schedule, fragment labelling and loss reduction shape source
retention and destination recognition.
The short-only state is worse than whole-word mixing by 7.328 destination WER
points and 6.238 source points, with substantially more deletions. Retaining
scheduled full views is therefore beneficial in this comparison; increasing
the proportion of short supervision does not monotonically improve transfer.

Two same-target controls restore neighboring pixels or supply predictions
informed by the complete video. Neither improves on FM in the Vietnamese
comparison. Their edit-error profiles show that improved visual coverage or
fewer substitutions can be offset by additional deletions, motivating assessment
of the complete output sequence.

The paper contributes three results:
\begin{itemize}
 \item A controlled study of source-only short-view continuation under matched
 update budgets, with unchanged whole-utterance recognition at deployment.
 \item A decoder-controlled comparison across two corpus pairs, including a
 two-seed Vietnamese transfer gain and joint source/destination assessment
 that identifies where the intervention does not improve recognition.
 \item Comparisons of temporal schedule, crop labelling, loss reduction and
 same-target context, supported by edit-error and exposure measurements.
\end{itemize}
Together, these results characterize when mixed temporal supervision improves
transfer, how retaining full utterances contributes, and the source-retention
cost of the resulting continuation strategy.
